% Zumdahl Chapter 7 test -- 20 MCQ + 2 FRQ. 50 min.
% Covers only the CED-relevant sections; skip-list material is not tested.
\documentclass{apchem-exam}

\apcunitword{Chapter}
\apcunit{7}
\apcunittitle{Atomic Structure and Periodicity}
\apcdoctitle{Chapter Test}
\apcsubtitle{Zumdahl \S7.1--7.2, 7.5, 7.8--7.9, 7.11--7.12 \textbullet\ 50 minutes}

\begin{document}
\apctitleblock

\begin{apcnote}[Scope]
This test covers only the CED-relevant sections. Nothing from \S7.3--7.4
(hydrogen spectrum, Bohr model), \S7.6--7.7 (quantum numbers, orbital
shapes), \S7.10, or \S7.13 appears.
\end{apcnote}

\examsection{I}{Multiple Choice}{20 questions \textbullet\ 30 minutes \textbullet\ 1 point each}{\nocalculator}

% ---- light and photons -----------------------------------------------------
\begin{mcq}
  Which relationship correctly connects the wavelength and frequency of
  light?
  \choices
    \correct{$\lambda\nu = c$}
    \choice{$\lambda + \nu = c$}
    \choice{$\lambda/\nu = c$}
    \choice{$\lambda\nu = h$}
\end{mcq}
\why{Speed is fixed, so wavelength and frequency are inversely related.}

\begin{mcq}
  Which type of electromagnetic radiation has photons of the greatest
  energy?
  \choices
    \choice{microwave}
    \choice{infrared}
    \choice{red visible light}
    \correct{ultraviolet}
\end{mcq}
\why{$E = hc/\lambda$ --- shortest wavelength, highest energy.}

\begin{mcq}
  A photon of \SI{300}{\nano\meter} light has how many times the energy of a
  photon of \SI{900}{\nano\meter} light?
  \choices
    \choice{$1/3$}
    \correct{3}
    \choice{9}
    \choice{the same}
\end{mcq}
\why{Energy is inversely proportional to wavelength.}
\distractor{C}{energy goes as $1/\lambda$, not $1/\lambda^2$}

\begin{mcq}
  Light below a metal's threshold frequency is shone on it at very high
  intensity. The result is
  \choices
    \choice{many electrons are ejected}
    \choice{a few electrons are ejected}
    \correct{no electrons are ejected}
    \choice{electrons are ejected only after a delay}
\end{mcq}
\why{Intensity supplies more photons, not more energetic ones.}
\distractor{D}{a wave model predicts this delay; it is never observed}

\begin{mcq}
  The photoelectric effect provides evidence that light
  \choices
    \choice{travels as a continuous wave only}
    \correct{is quantized into photons of energy $h\nu$}
    \choice{has mass}
    \choice{slows down in a vacuum}
\end{mcq}
\why{The threshold behaviour cannot be explained by a continuous wave.}

% ---- orbitals and shielding ------------------------------------------------
\begin{mcq}
  An atomic orbital is best described as
  \choices
    \choice{a circular path followed by an electron}
    \correct{a region of space where an electron is likely to be found}
    \choice{the surface of the nucleus}
    \choice{a fixed energy level containing all electrons}
\end{mcq}
\why{Orbitals are probability distributions, not trajectories.}

\begin{mcq}
  How many orbitals are in the $n = 3$ shell?
  \choices
    \choice{3}
    \choice{6}
    \correct{9}
    \choice{18}
\end{mcq}
\why{$1 + 3 + 5$ for $3s$, $3p$, $3d$.}
\distractor{D}{that is the maximum electron count}

\begin{mcq}
  In a polyelectronic atom, the 2$s$ orbital is lower in energy than the
  2$p$ because
  \choices
    \choice{2$s$ is farther from the nucleus on average}
    \correct{2$s$ penetrates closer to the nucleus and is shielded less}
    \choice{2$s$ holds more electrons}
    \choice{2$p$ orbitals are unstable}
\end{mcq}
\why{The small hump of 2$s$ density near the nucleus produces stronger net
attraction.}
\distractor{A}{true but incomplete --- and by itself would predict the
opposite}

\begin{mcq}
  In hydrogen, the 2$s$ and 2$p$ orbitals have
  \choices
    \correct{the same energy, because there is no shielding}
    \choice{different energies, because 2$s$ penetrates more}
    \choice{the same energy, because hydrogen has one proton}
    \choice{different energies, because of the Pauli principle}
\end{mcq}
\why{With only one electron there is no electron--electron repulsion to
break the degeneracy.}

\begin{mcq}
  A 2$p$ electron in fluorine experiences a larger effective nuclear charge
  than one in oxygen because fluorine has
  \choices
    \correct{one more proton with essentially the same shielding}
    \choice{one more electron shell}
    \choice{a larger atomic radius}
    \choice{fewer core electrons}
\end{mcq}
\why{$Z$ increases while $S$ stays essentially constant across a period.}

% ---- electron configurations -----------------------------------------------
\begin{mcq}
  The ground-state configuration of arsenic (33 electrons) is
  \choices
    \choice{$[\ce{Ar}]4s^2 4p^3$}
    \correct{$[\ce{Ar}]4s^2 3d^{10} 4p^3$}
    \choice{$[\ce{Ar}]4s^2 3d^{10} 4p^5$}
    \choice{$[\ce{Kr}]4s^2 3d^{10}$}
\end{mcq}
\why{$18 + 2 + 10 + 3 = 33$.}
\distractor{A}{omits the $3d$ subshell entirely}

\begin{mcq}
  The configuration of \ce{Fe^2+} is
  \choices
    \choice{$[\ce{Ar}]4s^2 3d^4$}
    \correct{$[\ce{Ar}]3d^6$}
    \choice{$[\ce{Ar}]4s^2 3d^6$}
    \choice{$[\ce{Ar}]4s^1 3d^5$}
\end{mcq}
\why{Transition metals lose 4$s$ electrons before 3$d$.}
\distractor{A}{removes $d$ electrons instead --- the standard error}

\begin{mcq}
  Which species is isoelectronic with argon?
  \choices
    \choice{\ce{Na+}}
    \choice{\ce{F-}}
    \correct{\ce{S^2-}}
    \choice{\ce{Mg^2+}}
\end{mcq}
\why{$16 + 2 = 18$ electrons.}

\begin{mcq}
  How many unpaired electrons are in a ground-state \ce{Fe^3+} ion?
  \choices
    \choice{3}
    \choice{4}
    \correct{5}
    \choice{0}
\end{mcq}
\why{$[\ce{Ar}]3d^5$ --- one electron in each of the five $d$ orbitals, by
Hund's rule.}

\begin{mcq}
  Which configuration represents an excited state of a neutral atom?
  \choices
    \choice{$1s^2 2s^2 2p^6$}
    \correct{$1s^2 2s^2 2p^5 3s^1$}
    \choice{$1s^2 2s^2 2p^6 3s^1$}
    \choice{$1s^2 2s^2 2p^3$}
\end{mcq}
\why{Ten electrons (neon) with 2$p$ short while 3$s$ is occupied --- aufbau
violated.}

\begin{mcq}
  The actual ground-state configuration of chromium is
  \choices
    \choice{$[\ce{Ar}]4s^2 3d^4$}
    \correct{$[\ce{Ar}]4s^1 3d^5$}
    \choice{$[\ce{Ar}]4s^2 3d^5$}
    \choice{$[\ce{Ar}]3d^6$}
\end{mcq}
\why{A half-filled $d$ subshell is slightly favoured, and $4s$/$3d$ are
close in energy.}

% ---- trends and PES --------------------------------------------------------
\begin{mcq}
  Which atom has the largest atomic radius?
  \choices
    \correct{K}
    \choice{Ca}
    \choice{Na}
    \choice{Mg}
\end{mcq}
\why{Radius grows down a group and shrinks across a period; K is farthest
down and left.}

\begin{mcq}
  The first ionization energy of oxygen is lower than that of nitrogen
  because
  \choices
    \choice{oxygen has a larger atomic radius}
    \correct{oxygen's fourth 2$p$ electron is paired, adding repulsion}
    \choice{oxygen has fewer protons}
    \choice{nitrogen has a filled 2$p$ subshell}
\end{mcq}
\why{Nitrogen's 2$p$ electrons are all unpaired; oxygen must pair one.}
\distractor{D}{nitrogen's 2$p$ is half-filled, not filled}

\begin{mcq}
  Successive ionization energies for an element are 578, 1817, 2745, and
  11{,}577 kJ/mol. The element is in group
  \choices
    \choice{1}
    \choice{2}
    \correct{13}
    \choice{14}
\end{mcq}
\why{The large jump comes after the third electron, so three valence
electrons (it is aluminium).}

\begin{mcq}
  In a photoelectron spectrum, the area under a peak is proportional to
  \choices
    \choice{the binding energy of those electrons}
    \correct{the number of electrons in that subshell}
    \choice{the atomic number}
    \choice{the energy of the incident photons}
\end{mcq}
\why{Position gives energy; area gives population.}

\examsection{II}{Free Response}{2 questions \textbullet\ 20 minutes}{\calculatorok}

% ---- FRQ 1 -----------------------------------------------------------------
\frq{short}{4}{Zumdahl \S7.1--7.2 \textbullet\ photons and the photoelectric effect}

A certain metal has a threshold frequency of \SI{8.5e14}{\per\second}.
Constants: $c = \SI{3.00e8}{\meter\per\second}$,
$h = \SI{6.626e-34}{\joule\second}$.

\begin{parts}
  \item Calculate the minimum photon energy required to eject an electron
        from this metal. \workspace[2cm]
        \keyonly{\begin{apcanswer}$E = h\nu_0 =
        (6.626\times10^{-34})(8.5\times10^{14}) =
        \SI{5.6e-19}{\joule}$\end{apcanswer}}
  \item Determine whether light of wavelength \SI{400.}{\nano\meter} will
        eject electrons. Show your reasoning. \workspace[2.2cm]
        \keyonly{\begin{apcanswer}$\nu = 3.00\times10^{8}/400.\times10^{-9}
        = \SI{7.50e14}{\per\second}$, which is below
        $8.5\times10^{14}$ --- so no electrons are
        ejected.\end{apcanswer}}
  \item Explain why increasing the intensity of that \SI{400.}{\nano\meter}
        light would not change your answer to (b).
        \answerlines[3]{Intensity determines how many photons arrive, not
        how much energy each carries. Ejection requires a single photon with
        at least the threshold energy, so supplying more insufficient
        photons accomplishes nothing.}
\end{parts}

\begin{rubric}
  \rpoint{1}{(a) \SI{5.6e-19}{\joule} with units}
  \rpoint{1}{(b) converts wavelength to frequency correctly}
  \rpoint{1}{(b) concludes no ejection, by comparison to the threshold}
  \rpoint{1}{(c) distinguishes photon energy from photon number ---
    an answer about ``total energy'' earns 0}
\end{rubric}

% ---- FRQ 2 -----------------------------------------------------------------
\frq{short}{4}{Zumdahl \S7.11--7.12 \textbullet\ configuration, PES, and trends}

The complete photoelectron spectrum of an element shows peaks at 126, 9.07,
5.31, and 0.74 MJ/mol with relative areas $2 : 2 : 6 : 2$.

\begin{parts}
  \item Determine the electron configuration and identify the element,
        citing your evidence. \answerlines[2]{Areas total 12 electrons:
        $1s^2 2s^2 2p^6 3s^2$ --- magnesium.}
  \item Explain why the peak at 5.31 MJ/mol lies at lower binding energy
        than the peak at 9.07, even though both correspond to electrons in
        the $n = 2$ shell. \answerlines[3]{The 9.07 peak is 2$s$ and the
        5.31 peak is 2$p$. The 2$s$ orbital penetrates closer to the
        nucleus, so it is shielded less and held more tightly; 2$p$
        experiences slightly more shielding and a smaller
        $Z_{\text{eff}}$.}
  \item Predict whether aluminium's 1$s$ binding energy is greater or less
        than 126 MJ/mol, and justify with Coulombic reasoning.
        \answerlines[3]{Greater --- aluminium has 13 protons to magnesium's
        12, and 1$s$ electrons are essentially unshielded, so they
        experience nearly the full increase in nuclear charge at
        comparable distance.}
\end{parts}

\begin{rubric}
  \rpoint{1}{(a) configuration matching all four areas AND magnesium}
  \rpoint{2}{(b) 1 pt assigns 2$s$ and 2$p$ correctly; 1 pt penetration or
    shielding argument --- ``2$p$ is farther out'' alone earns 0}
  \rpoint{1}{(c) greater, justified by nuclear charge on unshielded 1$s$;
    stability or octet language earns 0}
\end{rubric}

\examtotal

\end{document}
